\section{Equilibrio, elasticidad y eficiencia del mercado}
\label{sec:equilibrio-elasticidad-eficiencia}

En este bloque se juntan las piezas principales de la microeconomía. Hasta aquí ya vimos qué es la demanda, qué es la oferta, y cómo se comportan los mercados. Ahora se estudia cómo se determinan precios y cantidades en equilibrio, cómo responden los consumidores y productores a cambios, y cómo se mide el bienestar social.

\subsection{Equilibrio de mercado}

\begin{concepto}{Equilibrio}
El equilibrio de mercado es el punto donde la cantidad demandada y la cantidad ofrecida son iguales.
\end{concepto}

En equilibrio, hay un precio y una cantidad que vacían el mercado. No hay exceso ni escasez.

\begin{equation}
Q_d = Q_s
\end{equation}

\begin{ejemplo}
Si el precio de las naranjas es demasiado bajo, hay exceso de demanda: la gente quiere comprar más de lo que hay disponible. Si el precio es demasiado alto, hay exceso de oferta: hay más naranjas de las que se venden. El mercado se mueve hacia el punto de equilibrio.
\end{ejemplo}

\subsection{Exceso de oferta y exceso de demanda}

\begin{itemize}
    \item \textbf{Exceso de oferta:} el precio está por encima del equilibrio; $Q_s > Q_d$.
    \item \textbf{Exceso de demanda:} el precio está por debajo del equilibrio; $Q_d > Q_s$.
\end{itemize}

\begin{ejemplo}
Si se pone un precio artificialmente alto sobre un bien, los productores quieren vender más, pero los consumidores compran menos. Esa diferencia genera excedente de inventario.
\end{ejemplo}

\subsection{Estática comparativa}

La estática comparativa analiza cómo cambia el equilibrio cuando se mueve una curva de demanda o de oferta.

\begin{concepto}{Estática comparativa}
Es el análisis del nuevo equilibrio tras un cambio en alguna variable exógena.
\end{concepto}

\textbf{Pasos clave:}
\begin{enumerate}
    \item identificar si se mueve la demanda o la oferta,
    \item determinar si se desplaza a la derecha o a la izquierda,
    \item ver cómo cambia el precio y la cantidad de equilibrio.
\end{enumerate}

\begin{ejercicio}{Mercado del sushi}
El sushi se torna muy popular y, al mismo tiempo, hay una crisis pesquera que reduce la disponibilidad de salmón.

- La demanda del sushi aumenta.
- La oferta de salmón disminuye.
- El precio del sushi sube.
- La cantidad puede subir, bajar o quedar incierta, dependiendo de qué shock sea más fuerte.
\end{ejercicio}

\subsection{Elasticidad precio de la demanda}

La elasticidad mide la sensibilidad de la cantidad demandada a un cambio en el precio.

\begin{equation}
\eta = \left| \frac{\%\Delta Q_d}{\%\Delta P} \right|
\end{equation}

\begin{itemize}
    \item $|\eta| = 0$: perfectamente inelástica.
    \item $0 < |\eta| < 1$: inelástica.
    \item $|\eta| = 1$: unitaria.
    \item $|\eta| > 1$: elástica.
    \item $|\eta| = \infty$: perfectamente elástica.
\end{itemize}

\begin{concepto}{Demanda inelástica}
La cantidad demandada responde poco a cambios en el precio.
\end{concepto}

\begin{concepto}{Demanda elástica}
La cantidad demandada responde mucho a cambios en el precio.
\end{concepto}

\begin{ejemplo}
La insulina tiene demanda muy inelástica. El precio del pan también suele ser inelástico en comparación con bienes de lujo o viajes.
\end{ejemplo}

\subsection{Determinantes de la elasticidad de la demanda}

\begin{itemize}
    \item existencia de sustitutos,
    \item proporción del ingreso gastado,
    \item necesidad del bien,
    \item horizonte de tiempo,
    \item definición del mercado.
\end{itemize}

\begin{ejemplo}
La demanda de gasolina es relativamente inelástica en el corto plazo porque hay pocos sustitutos inmediatos. En el largo plazo, con más opciones de transporte, la demanda puede volverse más elástica.
\end{ejemplo}

\subsection{Elasticidad y ingreso total}

\begin{equation}
IT = P \times Q
\end{equation}

\begin{itemize}
    \item Si la demanda es inelástica, un aumento del precio aumenta el ingreso total.
    \item Si la demanda es elástica, un aumento del precio reduce el ingreso total.
\end{itemize}

\begin{ejemplo}
Una empresa que vende un bien muy necesario puede subir el precio y aun así aumentar sus ingresos, porque pierde pocos clientes. En bienes de lujo, subir el precio puede hacer caer mucho las ventas.
\end{ejemplo}

\subsection{Elasticidad precio de la oferta}

\begin{equation}
\varepsilon = \frac{\%\Delta Q_s}{\%\Delta P}
\end{equation}

La elasticidad de la oferta mide cuánto responde la cantidad ofrecida a cambios en el precio.

\begin{itemize}
    \item oferta inelástica: difícil ajustar la producción rápidamente,
    \item oferta elástica: fácil ajustar la producción.
\end{itemize}

\begin{ejemplo}
En el corto plazo, una empresa puede no poder producir mucho más porque sus fábricas ya están operando. En el largo plazo, puede invertir y ampliar la producción.
\end{ejemplo}

\subsection{Excedente del consumidor}

\begin{concepto}{Excedente del consumidor}
Es la diferencia entre lo que un consumidor está dispuesto a pagar y lo que realmente paga.
\end{concepto}

El excedente del consumidor mide el bienestar de los compradores.

\begin{ejemplo}
Si una persona está dispuesta a pagar $15.000 por un libro, pero lo compra en $10.000, su excedente es $5.000.
\end{ejemplo}

\subsection{Excedente del productor}

\begin{concepto}{Excedente del productor}
Es la diferencia entre el precio que recibe el vendedor y su costo de producción.
\end{concepto}

Mide el bienestar de los productores.

\begin{ejemplo}
Si un productor vende un producto a $20.000 y su costo de producirlo es $12.000, su excedente es $8.000.
\end{ejemplo}

\subsection{Excedente total}

\begin{equation}
ET = EC + EP
\end{equation}

El excedente total representa el bienestar social generado por una transacción o por un mercado. El equilibrio competitivo suele maximizar el excedente total, siempre que no existan fallas de mercado.

\begin{advertencia}{Importante}
Si se produce demasiado poco o demasiado mucho, se destruye valor social. El equilibrio competitivo es eficiente porque maximiza el bienestar total en condiciones estándar.
\end{advertencia}

\subsection{Ejercicio resuelto: equilibrio y elasticidad}

\begin{ejercicio}{Mercado de sopaipillas}
La demanda es:
\[P = 500 - 6Q_d\]
La oferta es:
\[P = 50 + 3Q_s\]

\textbf{1. Hallamos el equilibrio:}
\[500 - 6Q = 50 + 3Q\]
\[450 = 9Q\]
\[Q^* = 50\]
Luego:
\[P^* = 500 - 6(50) = 200\]

\textbf{2. Interpretación:}
En equilibrio, se venden 50 sopaipillas a un precio de 200.

Si se fija un precio mínimo superior al equilibrio, habrá exceso de oferta. Si se fija un precio máximo menor al equilibrio, habrá exceso de demanda.
\end{ejercicio}

\subsection{Resumen final del bloque}

El equilibrio de mercado es el punto donde se igualan oferta y demanda. La elasticidad mide la sensibilidad al precio. El bienestar económico se observa en los excedentes del consumidor y productor. En condiciones normales, el equilibrio competitivo maximiza el bienestar total.

\begin{resumen}{Puntos clave}
- El equilibrio es $Q_d = Q_s$.
- Exceso de oferta y demanda aparecen cuando el precio está fuera del equilibrio.
- La elasticidad mide sensibilidad.
- El mercado eficiente maximiza el excedente total.
\end{resumen}
